% Transcription — fonds Alexandre Grothendieck, shelfmark 76, batch 2 (pages 21-40).
% Established from the facsimile archives/batches/76/batch-02.pdf, cut from a
% copy of 76.pdf fetched through the project's relay
% (https://grothendieck-relay.onrender.com/source/76.pdf); PDF page k =
% archivists' page 20+k, no cover sheet. Per the skill transcribe-grothendieck.
% The folder is sixty-two pages in four batches (1-20, 21-40, 41-60, 61-62),
% transcribed in parallel by four passes. Batch 1 existed when this pass
% finished and its notation was followed (D_*, \mathfrak{P}^*, (Ens), \Delta_n,
% \mathbf{N}); it was not edited.
% Pass: Opus 5.5 (claude-opus-5-5), 2026-09-26 — first pass, unchecked against the pages by a human.
%
% Dating and title. The dating is the folder's own, copied verbatim from the
% catalogue; the group « Géométrie et topologie combinatoire » (69-89) holds
% it, and since the folder has a date of its own the group's range is not
% added. The title is the catalogue's, without its final full stop, as in
% batch 1.
%
% Pages skipped: 34, a blank cream sheet (it separates two runs). Pages
% summarised: none; nothing is typed and nothing is by another hand.
%
% Pages transcribed from his hand: 21-33, 35-40, in black ink with a few
% additions in blue felt-tip on pp. 21-22. Pp. 21-33 are written fair and
% read almost end to end; p. 35 (the head of a new run) and p. 36 are
% faster, p. 36's lower third is cancelled by three long diagonal strokes
% and given only as far as it reads, as a single \struck block.
%
% His pagination: none on the leaves. Numbered runs and labels: circled
% conditions (1) (from p. 20), (2°) and (2 bis) on p. 22, (a), (b), (c), (d),
% (a') on pp. 21-25; « Ex 1 » (p. 25), « Ex 2 » (p. 26), a « Lemme » and its
% « Dém » (pp. 27-28); points a), b) on p. 29; on the new run from p. 35,
% circled cases (0), (1) (p. 35), (1) n = 1, (2) n >= 2 (p. 37), (1), (2)
% again (p. 40), and « Th. 1 », « Th 2 » (p. 38), « Th. 3 » (p. 39),
% « Th 4 » (p. 40, announced only). Each circled label is set as (n) with a
% note. No « cf. p. N » cross-reference occurs.
%
% What the batch is about. Pp. 21-33 continue p. 20 of batch 1 on
% quasi-simplicial sets (presheaves on Delta, finite non-empty totally
% ordered sets with strictly increasing maps) and ordered sets. P. 21:
% condition (a), each D_*(d) = {delta <= d} is some P^*(I_d); the question
% how the total orders of the I_d can come from one order on D_0. P. 22:
% the relation x_0 < y_0 iff some z_1 in D_1 has x_0 = delta_0 z_1,
% y_0 = delta_1 z_1; condition (2 bis), surjectivity of
% D_2 -> D_1 x_{D_0} D_1, gives transitivity, (1) antisymmetry, and the
% induced order on I_d is the given one. Pp. 23-24: conversely an ordered
% set with (a) and an order on D_0 satisfying (b), (c) gives a
% quasi-simplicial set with (1); the two constructions are inverse; the
% opposite order gives another quasi-simplicial set with homeomorphic
% realisation. P. 25, Ex 1: flags of an ordered set X, condition (d)
% d = Sup(I_d) characterising simplicial complexes, (a'). Pp. 26-28, Ex 2:
% a group G with elements sigma_0..sigma_n, sigma_i sigma_j = sigma_j
% sigma_i for j >= i+2, D(H) = X/N_H for a G-set X ("drapeaux de type H");
% transitivity reduced to surjectivity of D(H u H') -> D(H) x_{D(H n H')}
% D(H'), a Lemme proving it when H n H' is a final segment of H and an
% initial segment of H', via G_{K u K'} = G_K G_{K'}. P. 29: programme
% a) (surjectivity in general), b) cellular decomposition. Pp. 30-33:
% presheaves D on finite subsets of N, the ordered set Theta(D) and its
% graded pieces Theta_n(D), realisation |D|, the semi-simplicial structure
% sigma^* : Theta_m(D) -> Theta_n(D), generalisation to an ordered set N
% and its flags Dr(N). Pp. 35-40, a new run headed by a boxed « Cartes
% cellulaires n-dim et groupes cartographiques »: n-admissible (« régulières »)
% filtrations X_0 c ... c X_n = X built by attaching closed n-balls along a
% finite map q_n : boundary -> X_{n-1}; morphisms (finite, compatible,
% étale on strata), categories F_n and F_n°; the n-cartographic group G_n
% (sigma_i^2 = 1, far-apart sigma's commute); functors E_n : F_n ->
% G_{n-1}-ens, E_n° : F_n° -> G_n-ens by induction (E_1(X_1) = boundary of
% the decomposed 1-skeleton, sigma_0 and sigma_1 the characteristic
% involutions); Th. 1 (E_n° fully faithful on the isotopy category, image
% = G_n-sets on which the sigma_i act freely), cascades R_n(X_n) -> ... ->
% R_0(X_n), Th. 2 (F_n -> (Casc)_n fully faithful), Th. 3 (the cascade of an
% X_n in F_n° is E_n°(X_n) modulo sigma_n, then (sigma_n, sigma_{n-1}), ...),
% and a construction psi_n : (Casc)_n -> F_n by « multicône », with Th 4
% (description and adjunction) announced.
%
% Notation, fixed once for the batch and aligned with batch 1: D_*, D_n,
% \underline{D}_* for the fibred category, \mathfrak{P}^*, \mathfrak{P}f,
% (Ens), \Delta, \Delta_n = [0,n], \mathbf{N}. On p. 21 he writes the faces
% \partial_i, from p. 22 \delta_i; kept as written. The new order on D_0 is
% \prec, and \prec_* the one read off the quasi-simplicial structure (p. 23);
% his cursive X on p. 25 is \mathfrak{X}. « ½simpl » is \tfrac{1}{2}\mathrm{simpl}
% (semi-simplicial, hand.md). His script F of the new run is \mathcal{F}_n;
% \widetilde{X}_i, \partial\widetilde{X}_i as he draws them; (Casc)_n.
%
% Hardest pages: 22 (a blue diagonal marginal note, mostly lost), 35 (the
% boxed title and the cancelled right-hand block), 36 (the diagonal left
% margin note and the cancelled case n = 1), 40 (the letter rho and
% « multicône »). Readings to check first: « isotypique » vs « isotopique »
% (pp. 38-39; the second occurrence reads « isotopique »); « passe » (p. 38);
% « (Casca) » (p. 38); the sigma index on p. 37 (« noté sigma_n »).

\documentclass[11pt,a4paper]{article}
% Le préambule commun à toutes les transcriptions du fonds Grothendieck.
%
% Il tient en une idée : **l'incertitude se compose, elle ne se résout pas.**
% Une transcription qui lisse un mot douteux a détruit la seule chose pour
% laquelle on la faisait — un lecteur doit pouvoir savoir, à chaque endroit,
% ce qui était sur le feuillet et ce qui a été ajouté. Les six macros qui
% suivent sont donc l'appareil critique, et non de la décoration.
%
% Les mêmes commandes sont reconnues par `scripts/render.mjs`, qui produit la
% vue de lecture du site. Une macro ajoutée ici doit l'être aussi là-bas,
% faute de quoi le rendu échouera bruyamment — ce qui est voulu : un rendu
% silencieusement faux est pire qu'un rendu absent.

\ProvidesPackage{grothendieck}[2026/08/08 Transcriptions du fonds Grothendieck]

\usepackage{amsmath,amssymb,amsthm}
\usepackage{mathrsfs}
\usepackage{stmaryrd}
% Les diagrammes commutatifs sont partout dans ce fonds : c'est la langue
% même de la géométrie algébrique de Grothendieck, et une transcription qui
% les rendrait en prose ne transcrirait rien.
\usepackage{tikz-cd}
% Français par défaut — c'est la langue des feuillets — mais l'anglais est
% chargé aussi : la traduction et le résumé sont des documents anglais, et la
% typographie française (espaces avant les deux-points) y serait une faute.
% Ces documents basculent par \selectlanguage{english} en tête de corps.
\usepackage[english,french]{babel}
\usepackage{fontspec}
\usepackage{geometry}
\geometry{a4paper,margin=25mm,marginparwidth=28mm}
\usepackage{marginnote}
\usepackage{xcolor}
% ulem, not soul. `soul`'s \st and \ul are built on a character-by-character
% scan that XeLaTeX's font machinery breaks: the document fails with an
% undefined control sequence rather than degrading. ulem does the same two jobs
% and is Unicode-engine safe. `normalem` stops it hijacking \emph, which the
% transcriptions use for Grothendieck's own emphasis.
\usepackage[normalem]{ulem}
\usepackage{hyperref}
\hypersetup{colorlinks=true,linkcolor=black,urlcolor=black,pdfborder={0 0 0}}

% --- Métadonnées du lot -----------------------------------------------------
% Reprises telles quelles par le script de rendu, qui les affiche en tête de
% la vue de lecture. Elles fixent ce que le fichier prétend couvrir : sans
% elles, une transcription flotte sans point d'attache dans un fonds de seize
% mille feuillets.

\newcommand{\folder}[1]{\gdef\grothFolder{#1}}
\newcommand{\batch}[1]{\gdef\grothBatch{#1}}
\newcommand{\leaves}[2]{\gdef\grothFirst{#1}\gdef\grothLast{#2}}
\newcommand{\foldertitle}[1]{\gdef\grothFoldertitle{#1}}
\gdef\grothFolder{?}\gdef\grothBatch{?}\gdef\grothFirst{?}\gdef\grothLast{?}\gdef\grothFoldertitle{}

% --- Le résumé --------------------------------------------------------------
%
% Il ouvre la lecture modernisée, sous le titre « Résumé », et il est composé
% en petit. Ce n'est pas un ornement : c'est le seul passage du document qui
% ne soit pas des mathématiques, et un lecteur doit pouvoir le reconnaître
% avant de l'avoir lu — puis le sauter s'il connaît déjà le sujet, ou s'y
% arrêter s'il ne le connaît pas. Le filet qui le suit marque la reprise du
% corps.

\newenvironment{resume}
  {\small\setlength{\parskip}{0.45em}}
  {\par\medskip\hrule\medskip}

% --- Mots-clés ---------------------------------------------------------------
%
% En anglais, en clôture du résumé de la lecture modernisée : le vocabulaire
% moderne sous lequel chercher ce que les feuillets construisent. Le manifeste
% du site (`npm run manifest`) les extrait pour étiqueter la cote entière —
% cette ligne est la source unique des tags, il n'y a pas de fichier à côté.
\newcommand{\keywords}[1]{\par\smallskip\noindent
  {\footnotesize\textsc{Keywords} --- \textit{#1}}\par}

% --- L'appareil critique ----------------------------------------------------

% Le numéro de feuillet, en marge. C'est l'ancre qui permet de régler un
% désaccord sur une formule en regardant *un* feuillet plutôt qu'une chemise
% de six cents. Le site s'en sert aussi pour tourner les pages du fac-similé
% quand on fait défiler la transcription.
\newcommand{\leaf}[1]{%
  \marginnote{\footnotesize\color{gray!70}\textsf{#1}}%
  \label{leaf:#1}\ignorespaces}

% Illisible : on ne devine pas. Un mot inventé qui se lit comme les autres est
% le pire résultat possible.
\newcommand{\ill}{\textcolor{red!60!black}{[\,\ldots\,]}}

% Lecture douteuse : le mot est donné, et signalé comme incertain.
\newcommand{\uncertain}[1]{\uline{#1}}

% Ajout de l'éditeur — une lettre manquante, un mot sous-entendu.
\newcommand{\add}[1]{\textcolor{blue!50!black}{[#1]}}

% Biffé par Grothendieck lui-même. Ce qu'il a rayé fait partie du document :
% une rature dit souvent où la pensée a changé de direction.
\newcommand{\struck}[1]{\sout{#1}}

% Note du transcripteur, sur l'état matériel du feuillet.
\newcommand{\note}[1]{\par\smallskip{\small\color{gray!60!black}\textit{[#1]}}\par\smallskip}

% Note marginale de Grothendieck, distincte du corps du texte.
\newcommand{\marginal}[1]{\par\smallskip\noindent\hspace{1em}%
  {\small\color{gray!60!black}#1}\par\smallskip}

% --- Titre ------------------------------------------------------------------

\AtBeginDocument{%
  \begin{center}
    {\large\bfseries Fonds Alexandre Grothendieck --- cote n\textsuperscript{o}~\grothFolder}\\[2pt]
    {\itshape\grothFoldertitle}\\[4pt]
    {\small feuillets \grothFirst--\grothLast\ (lot \grothBatch)}\\[2pt]
    {\footnotesize\color{gray!60!black}Transcription établie d'après le fac-similé numérisé
     par l'Université de Montpellier.\\
     Les lectures incertaines sont soulignées, les illisibles notées [\,\ldots\,],
     les ajouts entre crochets.}
  \end{center}
  \vspace{1em}}

\endinput


\folder{76}
\batch{2}
\pages{21}{40}
\dating{[vers 1977]}
\watermark{Édition de démonstration}
\foldertitle{n-cartes cellulaires : notes manuscrites (s.d.)}

\begin{document}

\page{21}
\note{la page continue la fin de la page 20 (réalisation topologique de
$D_*$ comme limite inductive de simplexes types)}
donc $|D_*|$ est aussi la réalisation topologique de l'ens
quasi-simplicial $D_*$.

L'ens ordonné $D_*$ a la propriété suivante

(a) $\forall\, d \in D_*$, l'ens \add{ordonné} $D_*(d) = \lbrace \delta
\in D_* \mid \delta \leqslant d \rbrace$ est isomorphe à un
$\mathfrak{P}^*(I_d)$, pour $I_d$ ens fini (déterminé à isom.\ près comme
l'ens des él.\ minimaux de $D_*(d)$, ou encore l'ens des éléments minimaux
de $D_*$ --- i.e.\ des él.\ de $D_0$ --- majorés par $d$)
\note{la lettre (a) est cerclée ; il renvoie ensuite aux conditions
cerclées (a), (b), (c), (d), données ici entre parenthèses}

Mais comment exprimer l'ordre total des $I_d$ \add{($d \in D_*$)} en
termes d'une structure \uncertain{q-simpl.}\ sur l'ens ordonné $D_*$ ?
[\uncertain{Souvent} elle sera induite par une structure d'ordre (pas
total) sur $D_0$, ens des él.\ minimaux de $D_*$.] La donnée de cet
ordre total, avec la condition que pour $d' < d$, l'ordre total de
$I_{d'}$ est induit par celui de $I_d$, \emph{équivaut} : à la donnée de
l'ens.\ quasi-simplicial $D_*$. On va examiner des conditions \add{sur
$D_*$} moyennant lesquelles les ordres totaux des $I_d$ ($d \in D_*$)
sont bien induits par un ordre unique sur $D_0$.
\note{un double trait vertical au crayon bleu, dans la marge gauche,
signale les trois dernières lignes de ce paragraphe}

Considérons

\[
  D_0 \;\overset{\partial_0}{\underset{\partial_1}{\leftleftarrows}}\; D_1
  \;\overset{\partial_0,\ \partial_1,\ \partial_2}{\Lleftarrow}\; D_2
\]
\note{deux flèches $\partial_0, \partial_1$ de $D_1$ vers $D_0$, trois
flèches $\partial_0, \partial_1, \partial_2$ de $D_2$ vers $D_1$}

Considérons sur $D_0$ la relation

\page{22}
\[
  x_0 \prec y_0 \;\overset{\text{déf}}{\Longleftrightarrow}\; \exists\, z_1
  \in D_1 \text{ t.q. } x_0 = \delta_0(z_1),\ y_0 = \delta_1(z_1)
\]
\note{il note $\prec$ (un « $<$ » arrondi) la relation nouvelle sur
$D_0$, et $<$ l'ordre de $D_*$ ; les bords $\partial_i$ de la page 21
deviennent ici $\delta_i$}
la condition \add{\uncertain{et} $x_0 < x_0$ dans l'« \ill{} » \ill{}}
\struck{\ill{}}\note{l'insertion est au crayon bleu, de même que le
crochet qui ouvre la définition}

\struck{\uncertain{Est-t-on} \uncertain{transitivité} \ill{} cette relation
$\exists$ \ill{}}

(2°) La relation $\prec$ \struck{\uncertain{est bien}} sur $D_0$ est
transitive. ? Oui si on suppose satisfaite la condition suivante
\note{le « 2° » est cerclé}

(2 bis) L'application canonique
\[
  D_2 \xrightarrow{(\delta_2, \delta_0)} (D_1, \delta_1) \times_{D_0}
  (D_1, \delta_0)
\]
est surjective,
\note{« 2 bis » est cerclé, au-dessus d'un autre numéro cerclé, noirci}

\marginal{\ill{} \ill{} une fonction « dimension » $\delta$ et on exige
(pour $x_0 \prec y_0$) $\delta(x_0) < \delta(y_0)$ \ill{} \ill{}}
\note{note en diagonale au crayon bleu dans le coin supérieur gauche,
encadrée, reliée par un trait à « $x_0 \prec y_0$ » en tête de page}

En effet si $x_0 \prec y_0$ et $y_0 \prec z_0$, i.e.
\[
  x_0 = \delta_0(z_1),\ y_0 = \delta_1(z_1),\ \struck{y_0} = \delta_0(z'_1),\
  z_0 = \delta_1(z'_1),
\]
$\exists\, u_2 \in D_2$ avec $z_1 = \delta_2(u_2)$, $z'_1 = \delta_0(u_2)$,
posons alors $z''_1 = \delta_1(u_2)$, on a
\[
  \delta_0(z''_1) = \delta_0\delta_1(u_2) = \delta_0\delta_2 u_2 =
  \delta_0(z_1) = x_0, \quad \delta_1(z''_1) = \delta_1\delta_1(u_2) =
  \delta_1\delta_0(u_2) = \delta_1(z'_1) = z_0,
\]
donc $x_0 \prec z_0$, cqfd. \struck{Je dis que la relation de}
\note{les identités $\delta_0\delta_1 = \delta_0\delta_2$ et
$\delta_1\delta_1 = \delta_1\delta_0$ sont écrites telles quelles}

\add{\uncertain{préordre}} Ainsi, posant
\[
  x_0 \preccurlyeq y_0 \;\overset{\text{déf}}{\Longleftrightarrow}\;
  x_0 = y_0 \text{ ou } x_0 \prec y_0
\]
(moyennant (2°)), on trouve une relation de \struck{\uncertain{pré}}ordre
dans $D_0$, c'est une relation d'ordre, car $x_0 \prec y_0$ et $y_0 \prec
x_0$ sont \uncertain{incompatibles}
\marginal{sinon $x_0 \prec x_0$ i.e.\ $\exists\, z_1$ avec $\delta_0(z_1) =
\delta_1(z_1) = x_0$ impossible par (1).}
\note{note de la marge gauche, reliée par une accolade ; le « 1 » est
cerclé et renvoie à la condition (1) de la page 20}
Ceci posé, si $d \in D_*$, on voit que la relation d'ordre sur
\[
  I_d = \lbrace x \in D_0 \mid x < d \rbrace
\]
est induite par la relation précédente. En effet, si $x, y \in D_0$, $x, y
< d$, si on a $x < y$ dans \struck{\ill{}} $I_d$, on a $x \prec y$, i.e.\
\add{l'inclusion} $I_d \hookrightarrow D_0$ est strictement croissante,
d'où on déduit facilement (l'ordre de $I_d$ étant total) \add{et $D_0$
étant ordonné (\uncertain{voir paragraphe précédent})} que l'ordre de
$I_d$ est induit par celui de $D_0$.

\marginal{Cet ordre sur $D_0$ satisfait la condition (b) Si $x_0 \prec
y_0$ ($x_0, y_0 \in D_0$) \ill{} \struck{\ill{}} $z \in D_*$, $x_0 < z$,
$y_0 < z$.}
\note{suite de la note de la marge gauche ; le « b » est cerclé ; un mot
biffé, avant « $z \in D_*$ », est peut-être un $\exists$}

\page{23}
Inversement, soit $D_*$ un ens.\ ordonné satisfaisant (a) ci-dessus.
Donnons-nous sur $D_0$ (ens des él.\ minimaux de $D_*$) une structure
d'ordre (pas seulement de préordre) \add{notée $x \preccurlyeq y$},
\struck{induisant \ill{}} satisfaisant (b) ci-dessus et de plus
\struck{\uncertain{vérifiant que}} (c) \add{Elle induit sur} $I_d$
\struck{($d \in D_*$)} $= \lbrace x \in D_0 \mid x \leqslant d \rbrace$
un ordre total \struck{qui} \add{ce dernier} sera évidemment fonctoriel
\struck{\uncertain{pour ordonné} (pour l'ordre induit par)} en $d$, de
sorte que $D_*$ définit un foncteur covariant $\underline{D}_* \to
\Delta$, qui est \emph{fibrant}, à fibres discrètes --- donc associé à un
foncteur $D_* : \Delta^\circ \to (\mathrm{Ens})$, et $D_0 =
D_*(\Delta_0)$ (plus général\supplied{ement} $D_n = D_*(\Delta_n)$ est l'ens
des $d \in D_*$ tels que $\operatorname{card}(I_d) = n+1$, i.e.\ $I_d
\simeq \Delta_n$), et $D_*$ satisfait évidemment (1).
\note{les lettres des conditions (a), (b), (c) et le chiffre (1) sont
cerclés}

\note{à gauche du paragraphe suivant, un schéma de deux flèches
verticales opposées (vers le bas, vers le haut) entre les deux classes
d'objets :}
\begin{quote}
Ens.\ quasi-simpliciaux \add{$D_*$} satisfaisant (1) et \uncertain{(2)}

$\downarrow \qquad \uparrow$

Ens ordonnés \add{$D_*$} satisfaisant (a) et munis d'une relation d'ordre
sur $D_0 =$ ens des él.\ minimaux de $D_*$, satisfaisant les conditions
(b) et (c)
\end{quote}

Soient donc $x_0, y_0 \in D_0$ tels que $x_0 \prec_* y_0$ i.e.\
$\exists\, z_1 \in D_1$, $x_0 = \delta_0(z_1)$, $y_0 = \delta_1(z_1)$ avec
$z_1 \in D_1$. Donc $I_{z_1} = \lbrace x_0, y_0 \rbrace$, et par
définition de $\delta_0, \delta_1$, $x_0$ est le 1er
él.\ de $I_{z_1}$ pour l'ordre \add{(total)} induit par $D_0$, $y_0$ le
deuxième, i.e.\ on a $x_0 \prec y_0$. Inversement, supposons $x_0 \prec
y_0$, \struck{donc} $\exists\, z \in D_*$ tel que $x_0 < z$, $y_0 < z$
--- on aura $z \in D_n$ avec $n \geqslant 1$, $x_0, y_0$ él.\ distincts de
$I_z$, donc aussi $x_0 \prec y_0$ pour
\note{il distingue ici $\prec_*$ (l'astérisque sous le signe), la
relation définie par la structure quasi-simpliciale comme à la page 22,
de $\prec$, l'ordre donné sur $D_0$}

\page{24}
l'ordre induit par $(D_0, \prec)$, \struck{ce qui} correspond à une
\struck{\ill{}} \add{morphisme} $\sigma : \Delta_1 \to \Delta_n$, telle
que, posant $z_1 = \sigma^*(u)$, on ait $x_0 = \delta_0(z_1)$, $y_0 =
\delta_0(z_1)$, i.e.\ on a $x_0 \prec_* y_0$.
\note{le second $\delta_0$ est écrit tel quel ; on attend $\delta_1$,
comme à la page 23}
Ainsi l'ordre $\prec$ donné : l'ordre sur $D_0$ (satisfaisant (b)) se
récupère : partir de la donnée quasi-simpliciale $D_*$ qu'il définit,
comme étant $\prec_*$.

\emph{Rmq} Supposons donné l'ens ordonné $D_*$, \add{satisfaisant (a),}
et supposons qu'il existe un ordre $\omega$ sur $D_0$ satisfaisant (b).
Alors l'ordre opposé $\omega^\circ$ satisfait également (b), mais donne
lieu à une autre loi quasi-simpliciale sur la famille des
\[
  D_n = \lbrace d \in D_* \mid \operatorname{card} I_d = n+1 \rbrace
\]
(sauf si $D_1 = 0$ i.e.\ l'ordre de $D_*$ est discret). Néanmoins, comme
les \add{réalisations top.\ des} deux ens.\ quasi-simpliciaux obtenus
sont \uncertain{can.}\ homéom.\ à la réalisation top $|D_*|$ de l'ens
ordonné $D_*$, (qui, elle, ne dépend pas du choix d'un ordre sur $D_0$)
on trouve \struck{que les 2} \add{un homéom.\ can} \struck{ens
quasi-simpliciaux}
\[
  |D_*^\omega| \simeq |D_*^{\omega'}|
\]
\note{il écrit $\omega'$ dans la formule finale pour l'ordre opposé noté
$\omega^\circ$ plus haut}

\page{25}
\emph{Ex 1} Soit $\mathfrak{X}$ un ens ordonné. Soit $D_*$ l'ens des
parties finies de $\mathfrak{X}$, totalement ordonnées par l'ordre induit
(« drapeaux de $\mathfrak{X}$ ») ; on ordonne $D_*$ par inclusion. Alors
$D_0$ (ens des points) s'identifie à $\mathfrak{X}$, il est
\struck{ordonné} donc muni d'une relation d'ordre, et (a) et (b) sont
évidemment satisfaits. D'où un ens.\ quasi-simplicial $D_*$, et on a des
homéom.\ \uncertain{can.}
\[
  |D_*| \simeq |D_*| \simeq |\mathfrak{X}|
\]
\note{la lettre qu'il écrit $\mathfrak{X}$ est un X cursif ; dans la
formule, le premier $D_*$ est l'ens.\ quasi-simplicial, le second l'ens
ordonné, qu'il distingue par la forme du D}

En fait, les $(D_*, \prec)$ obtenus sont ceux pour lesquels on a, en plus
de (a), (b), \add{(c)} la condition

(d) $\forall\, d \in D_*$, on a $d = \operatorname{Sup}(I_d)$

i.e.\ l'ens ordonné $D_*$ est l'ens des \emph{simplexes} \add{non vides}
d'une \emph{multiplicité simpliciale}, \struck{où l'ens quasi-simplicial
des} \struck{$D_*$}. Donc il faut prendre les multiplicités simpliciales
avec une structure d'ordre sur l'ens $\mathfrak{X} = D_0$ des sommets,
induisant un ordre total sur chaque simplexe, et \struck{tel que}
\add{(parmi les parties finies non vides)} (sur les simplexes
seulement…

[Notons aussi que \struck{moyennant (c) + (d) équivaut} \add{(a) + (d)
(i.e.\ $D_*$ correspond} à une multiplicité simpliciale) équivaut :
(a') + (d), avec

(a') Dans $D_*$ les Sup \add{(finis)} majorés existent…]
\note{les lettres (a), (b), (c), (d), (a') sont cerclées}

\page{26}
\emph{Ex 2} Soit $G$ un groupe, et soit $(\sigma_i)_{0 \leqslant i
\leqslant n}$ une système d'éléments de $G$, tels que $\sigma_i,
\sigma_j$ commutent pour \struck{$i < j$} $n \geqslant j \geqslant i+2
\geqslant i \geqslant 0$. \struck{\ill{} \ill{} \ill{} Groupes, et} pour
toute partie finie $H$ de $\Delta_n = [0,n]$, soit \struck{\ill{}}
$N_H =$ ss-groupe de $G$ engendré par $\lbrace \sigma_i \mid i \in
\Delta_n \setminus H \rbrace$. Donc $H \subset H' \Rightarrow N_{H'}
\subset N_H$. Soit $X$ un $G$-ens et soit
\[
  D(H) = X / N_H \qquad \text{(drapeaux de type } H\text{)}
\]
donc
\[
  H \subset H' \Longrightarrow i^*_{H,H'} : D(H') \to D(H),
\]
$H \mapsto D(H)$ foncteur contravariant
\[
  \underline{\mathfrak{P}(\Delta_n)}^\circ \longrightarrow (\mathrm{Ens})
\]
Soit
\[
  D_* = \coprod_{H \in \mathfrak{P}(\Delta_n)} D(H)
\]
relation d'ordre : $d \in D(H)$, $d' \in D(H')$,
\[
  d \leqslant d' \;\overset{\text{déf}}{\Longleftrightarrow}\; H \subset
  H' \text{ et } d = i^*_{H,H'}(d').
\]
(\uncertain{on voit}, $D_*$ satisfait la condition (a) : si $d \in D_*$,
$D_*(d) \simeq \mathfrak{P}(I_d)$ (où $I_d \simeq H$ si $d \in D(H)$)).

Les él.\ minimaux sont ceux de
\[
  D_0 = \coprod_{0 \leqslant i \leqslant n}
  \underbrace{D(\lbrace i \rbrace)}_{X / (\sigma_0 \cdots \widehat{\sigma}_i
  \cdots \sigma_n)}
\]
On va introduire sur $D_0$ une relation d'ordre (pour l'instant on n'a pas
utilisé la structure d'ordre total

\page{27}
sur l'ens d'indices $\Delta_n$ pour les $\sigma_i \in G$]

Si \struck{\ill{}} $x_i \in D(\lbrace i \rbrace)$, $y_j \in D(\lbrace j
\rbrace)$, on pose $x_i \prec y_j$ ssi $i < j$ et $\exists\, z \in D_*$
tel que \struck{\ill{}} $x_i, y_j \leqslant z$.
\struck{Cette relation d'ordre satisfait (b) (et (c) (\uncertain{par
définition}) évidemment)}

Cette relation est-elle transitive ? Supposons $i < j < k$, $x_i \in
D(\lbrace i \rbrace)$, $y_j \in D(\lbrace j \rbrace)$, $z_k \in D(\lbrace
k \rbrace)$, \struck{\ill{}} $x_i, y_j$ provenant de $u_{ij} \in
D(\lbrace i, j \rbrace)$, $y_j, z_k$ provenant de $v_{jk} \in D(\lbrace
j, k \rbrace)$. Donc $(u_{ij}, v_{jk})$ est un élément de
\[
  D(\lbrace i, j \rbrace) \times_{D(\lbrace j \rbrace)} D(\lbrace j, k
  \rbrace)
\]
on serait content qu'il provienne d'un élément $w_{ijk} \in D(\lbrace i,
j, k \rbrace)$, car ce dernier définirait un élément \struck{\ill{}}
$x_{ik} \in D(\lbrace i, k \rbrace)$ majorant $x_i$ et
\uncertain{$z_k$}. Donc \struck{\ill{}} il suffit que
\[
  D(\lbrace i, j, k \rbrace) \xrightarrow{\text{épi}} D(\lbrace i, j
  \rbrace) \times_{D(\lbrace j \rbrace)} D(\lbrace j, k \rbrace).
\]
\note{« épi » est entouré, au-dessus de la flèche}
Or je dis plus généralement ceci.

\emph{Lemme} Soient $H, H' \subset \Delta_n$ tels que $H \cap H' \neq
\emptyset$ et $H \cap H'$ soit un \emph{segment} final de l'ens tot.\
ordonné $H$, un \emph{segment} initial de l'ens tot ordonné $H'$.
Considérons l'application canonique
\[
  D(H \cup H') \longrightarrow D(H) \times_{D(H \cap
  H')} D(H')
\]
Celle-ci est \emph{surjective}.
\note{l'indice du produit fibré, d'abord écrit puis noirci, est récrit
en dessous : $D(H \cap H')$}

\page{28}
\emph{Dém} Soit \add{pour $K \subset \Delta_n$,} $G_K =$ ss-groupe de $G$
engendré par $\lbrace \sigma_i \mid i \in K \rbrace$. Soient
\struck{$K = \complement H$, $K' = \complement H'$} $K = \Delta_n
\setminus H$, $K' = \Delta_n \setminus H'$, donc $K \cup K' = \Delta_n -
H \cap H'$, la \struck{\ill{}} \add{flèche} envisagée est
\[
  X / G_{K \cap K'} \xrightarrow{\;p\;} X/G_K \times_{X/G_{K \cup K'}}
  X/G_{K'}
\]
\note{sous $X/G_{K \cap K'}$, une flèche verticale monte de $X$, et une
flèche oblique $p'$ va de $X$ au produit fibré}
Il suffit de prouver que $p'$ est surjectif. Je dis que l'on a
\[
  G_{K \cup K'} = G_K \cdot G_{K'}
\]
(d'où \struck{\ill{}} résultat s'ensuit aussitôt).

Or \struck{\ill{}} \add{soit} $i$ le dernier él.\ \struck{de $K$}
\add{de $H$ \struck{\ill{}} de $H \cap H'$} \struck{\ill{}}
\add{$i \in H \cap H'$}, $i \notin K \cup K'$, donc
\struck{$K_0 = K \cap [0, i-1] \subset K'$}

\marginal{OPS $H \cap H' \neq H, H'$}
\[
  K \cup K' = (K \cup K') \cap \underbrace{\complement \lbrace i
  \rbrace}_{[0, i-1] \sqcup [i+1, n]} = \overbrace{(K \cup K') \cap [0,
  i-1]}^{K'_0} \sqcup \underbrace{(K \cup K') \cap [i+1, n]}_{K_0}
\]
le dernier sommande est $\subset K$ car $[i+1, n] \subset K$, le premier
est $\subset K'$, car pour les $j \leqslant i-1$, \struck{\ill{}} si $j
\notin H \cap H'$ (i.e.\ $j \in K \cup K'$) \struck{\ill{}} $j \notin H'$
\struck{\ill{}}, car $H \cap H'$ est \emph{segment} initial de $H'$,
d'extrémité $i$, donc si $j \in H'$ \uncertain{on aurait} $j \in H$ donc
$j \in H \cap H'$) donc $j \notin H'$.

\struck{Donc $G_{K \cup K'} \subset G_K \cdot G_{K'}$, d'où cqfd}

Or $G_{K_0}$ et $G_{K'_0}$ \emph{commutent} donc
\[
  G_{K \cup K'} \subset \underbrace{G_{K_0} \cdot G_{K'_0}}_{\text{ss-gr}}
  \subset G_K \cdot G_{K'} \quad \text{cqfd.}
\]

\marginal{On a utilisé : a) le dernier él.\ de $H$ est dans $H'$ ; b) $H
\cap H'$ est un segment initial de $H'$}
\note{note oblique dans la marge gauche, sous un petit croquis de deux
segments qui se chevauchent, le point $i$ marqué à la fin du premier}

\page{29}
\emph{Examinons, en général}

a) Conditions pour que (2 bis) soit satisfait [en termes d'ens.\ ordonnés
$(D_*, \prec$ sur $D_0)$] et plus gén.\ que l'on ait surjectivité des
\[
  D(\Delta_n) \longrightarrow D(H) \times_{D(H \cap H')} D(H')
\]
pour $\Delta_n = H \cup H'$, $H \cap H'$ segment final de $H$, initial de
$H'$.
\note{« 2 bis » est cerclé ; la source de la flèche, $D(\Delta_n)$, est
écrite sur une première lettre surchargée}

b) La décomposition cellulaire de $D_*$ associée à l'ordre sur $D_0$.
Essayons \struck{$D_*$, \ill{}} $(D_*, <, \prec)$ de
\uncertain{reconstituer} à l'aide de la structure d'espace topologique
\struck{cellulaire} avec décomposition en « cellules » (et coniques…)
et structures coniques…
\note{un double trait vertical au crayon, dans la marge, en face du
début de b) ; le reste de la page est blanc}

\page{30}
Supposons donné, pour toute partie finie $A$ \add{$\neq \emptyset$} de
$\mathbf{N}$, un ens $D(A)$ (« drapeaux de type $A$ ») et pour toute
inclusion $A \subset B$, une application « restriction » $D(B) \to D(A)$,
avec \uncertain{conditions} d'identité et de transitivité --- ce qui
signifie qu'on a un \struck{foncteur} \add{préfaisceau d'ens}
\struck{de la} sur la cat.\ \add{$C$} associée à l'ens ordonné des
parties finies de $\mathbf{N}$, i.e.\ $D \in \mathrm{Ob}\, \widehat{C}$.

\struck{\uncertain{Donnons}/Soit d'autre part $\Delta_n = [0,n] \subset
\mathbf{N}$ (pour $n \in \mathbf{N}$), et soit $\sigma : \Delta_n
\hookrightarrow \Delta_m$} Soit
\[
  \Theta(D) = \coprod_{\substack{A \subset \mathbf{N} \\ \text{et } A
  \text{ fini}}} D(A)
\]
\[
  \Theta_n(D) = \coprod_{\substack{A \subset \mathbf{N} \\
  \operatorname{card} A = n+1}} D(A) \qquad n \in \mathbf{N}
\]
donc
\[
  \Theta(D) = \coprod_{n \in \mathbf{N}} \Theta_n(D)
\]

Dans $\Theta(D)$ on a une relation d'ordre, en posant
\[
  \underset{D(A)}{d} \;\leqslant\; \underset{D(A')}{d'}
  \;\overset{\text{déf}}{\Longleftrightarrow}\; A \subset A' \text{ et } d =
  i^*_{A,A'}(d')
\]
\note{sous $d$ et $d'$, un signe d'appartenance couché renvoie à $D(A)$
et $D(A')$}

[On peut dire que $D \in \mathrm{Ob}\, \widehat{C}$ définit une catégorie
fibrée sur $C$ : fibres discrètes, et cette catégorie est associée à un
ens ordonné, qui est $\Theta(D)$ ; on a une application croissante
canonique
\[
  \Theta(D) \longrightarrow \mathfrak{P}_f(\mathbf{N})
\]

\page{31}
Nous nous intéressons notamment au cas où les $D(A)$ sont finis, et
vides sauf un ens fini d'entre eux, i.e.\ $\exists\, N \in \mathbf{N}$
tel que
\[
  A \not\subset [0, N] \Longrightarrow D(A) = \emptyset.
\]
Je remarque que $\Theta(D)$ est fini. La réalisation topologique
\add{$|\Theta(D)|$} de $\Theta(D)$ \struck{notée} est aussi notée $|D|$ et
s'appelle réalisation topologique de $D$. C'est un espace à
triangulation finie, \struck{\ill{}} et $\Theta(D)$ en tant qu'ens ordonné
se réalise comme ensemble de parties de $|D|$.

On va donner une deuxième description de $|D|$, en définissant
certaines opérations simpliciales entre les $\Theta_n(D)$. Posons
$\Delta_n = [0, n] \subset \mathbf{N}$, soit
\[
  \sigma : \Delta_n \hookrightarrow \Delta_m
\]
une application strictement croissante. Si $A \in
\mathfrak{P}f_{m+1}(\mathbf{N})$, on a une unique bij.\ croissante
\[
  \Delta_m \xrightarrow[\sim]{\;u_A\;} A
\]
d'où une partie finie de \add{card $n+1$} \struck{\ill{}}, $B =
u_A\sigma(\Delta_n)$, de $A$, et on \uncertain{considère}
\[
  \sigma^* : \lbrace D(A) \to D(B) \hookrightarrow \Theta_n(D) \qquad (B =
  u_A \sigma(\Delta_n))
\]
\note{$\mathfrak{P}f_{m+1}(\mathbf{N})$ : les parties de $\mathbf{N}$ à
$m+1$ éléments ; l'accolade ouvrante devant $D(A)$ est de lui, elle
annonce la définition de $\sigma^*$ par morceaux, $A$ variable}

\page{32}
donné par $\sigma^* = i^*_{B,A}$. Faisant varier $A$ dans
$\mathfrak{P}f_{m+1}(\mathbf{N})$, on trouve
\[
  \sigma^* : \Theta_m(D) \longrightarrow \Theta_n(D)
\]
avec transitivité évidente. Donc les $\Theta_n(D)$ ($n \in \mathbf{N}$)
forment un ens.\ quasi-simplicial (op.\ simpliciales définies par
applications \emph{str\supplied{ictemen}t croissantes}) i.e.\ définissent un
foncteur
\[
  \underbrace{(\tfrac{1}{2}\mathrm{simpl}^*)^\circ}_{\substack{\text{ens.\
  finis tot}^{\text{t}} \\ \text{ordonnés, avec} \\ \text{appl.\ str.\
  cr.}}} \longrightarrow (\mathrm{Ens})
\]
\note{« $\frac{1}{2}$simpl » : son abréviation pour « semi-simplicial » ;
la légende est écrite sous le symbole, sans accolade}

[Tout ceci marche en remplaçant $\mathbf{N}$ par un ens ordonné
quelconque \add{$N$}, $\mathfrak{P}f(\mathbf{N})$ par l'ens
\add{$\mathrm{Dr}(N)$} des parties finies tot\supplied{alemen}t ordonnées de
$N$, i.e.\ l'ens des drapeaux de $N$, qui est aussi l'ens ordonné associé
\add{à la cat.\ fibrée définie par le} \struck{au} foncteur
\[
  (\tfrac{1}{2}\mathrm{simpl}^*) \xrightarrow{\;\varphi_N\;} \mathrm{Ens}
\]
donné par
\[
  \varphi_N(\Delta_n) = \mathrm{Mon.croiss}(\Delta_n, N)
\]

\begin{tikzcd}
  \mathrm{Dr}(N) \arrow[d, "\text{appl.\ str.\ croissante}"] \\
  \tfrac{1}{2}\mathrm{simpl}
\end{tikzcd}


\page{33}
On a donc, \struck{des applications strictement croissantes}
\struck{\add{pour un \ill{} foncteur}}
\[
  D : \underline{\mathrm{Dr}(N)}^\circ \longrightarrow \mathrm{Ens},
\]
définissant une catégorie fibrée sur $\underline{\mathrm{Dr}(N)}$ associée
à un ens ordonné $\Theta(D)$, une appl.\ strictement croissante $\Theta(D)
\to \mathrm{Dr}(N)$, d'où appl.\ composée

\begin{tikzcd}
  \Theta(D) \arrow[d] \\
  \mathrm{Dr}(N) \arrow[d] \\
  (\tfrac{1}{2}\mathrm{simpl})^*
\end{tikzcd}

\note{la page s'arrête là, aux deux tiers vides}

\page{35}
\marginal{\struck{\uncertain{Divers}} \uncertain{Cartes} \uncertain{cellulaires}
$n$-dim et groupes cartographiques. \uncertain{fermées}}
\note{encadré de sa main en haut à droite de la page, souligné ; le
dernier mot, écrit sous le trait, est douteux. La page ouvre une nouvelle
suite, sans numéro de sa main}

\[
  X_{-1} = \emptyset \subset X_0 \subset X_1 \subset \cdots \subset
  X_{n-1} \subset X_n = X \qquad \text{« filtration } n\text{-admissible »}
\]
(\add{pour} $0 \leqslant i \leqslant n$)

a) \struck{Déc} $(X_i, X_{i-1})$ \add{$= \widetilde{X}_i$} est
\struck{une réunion disjointe} \add{somme top.} de boules de dim $i$, dont
la partie manquée $\Xi_i$ (image inverse de $X_{i-1}$ par $\widetilde{X}_i
\to X_i$) est le bord

\begin{tikzcd}
  \widetilde{X}_i \arrow[d, "p_i"'] & \supset & \partial\widetilde{X}_i
  \arrow[d, "q_{i-1}"] \\
  X_i & \supset & X_{i-1}
\end{tikzcd}

\note{il note $\operatorname{Déc}(X_i, X_{i-1}) = \widetilde{X}_i$ la
« décomposition » de $X_i$ le long de $X_{i-1}$, comme en batch 1 ; la
lettre qu'il écrit pour la partie manquée ressemble à un $\Xi$}

($\Rightarrow$ $X_i - X_{i-1}$ est une \struck{\ill{}} somme top.\ de
boules ouvertes (de dim $i$)), enfin $\widetilde{X}_i \to X_i$ est
morphisme fini.

b) \add{(pour $1 \leqslant i \leqslant n$)} \struck{la filtration} sur
$\partial\widetilde{X}_i$ image inverse de la filtration de
\[
  \emptyset \subset X_0 \subset X_1 \subset \cdots \subset X_{i-1}
\]
est \supplied{$(i-1)$-}admissible (NB le cas $i = 1$ est trivial), et les
applications $q_i^{-1}(X_j) - q_i^{-1}(X_{j-1}) \to X_j - X_{j-1}$ ($0
\leqslant j \leqslant i-1$) sont étales.

Ceci donne une définition récurrente de « filtration $n$-admissible »
pour $n \geqslant 0$
\struck{filtration $n$-admissible $\Leftrightarrow$
($\emptyset \subset X_0 \subset \cdots \subset X_{n-1}$) est
$(n-1)$-admissible, $\widetilde{X}_n$ \ill{} de $n$-boules fermées…}
\note{ce passage, dans la colonne de droite, est barré de plusieurs
traits obliques}

(0) Filtration $0$-admissible : $\emptyset \subset X_0 = X$ : $X$ discret
(néc.\ et \emph{suffisant} \uncertain{si déf})

(1) Filtration $1$-admissible : $\emptyset \subset X_0 \subset X_1$ :
$X_0$ discret, $\operatorname{Déc}(X_1, X_0) = \widetilde{X}_1$ somme top.\
de segments, dont partie manquée est $\partial\widetilde{X}_1$, et
$\widetilde{X}_1 \to X_1$ propre $\Leftrightarrow$ $(X_0, X_1)$
\uncertain{1-complexe} top

…
\note{les chiffres (0) et (1) sont cerclés}

Filtration $n$-admissible : se déduisent de filtration $(n-1)$-admissible
\[
  \emptyset \subset X_0 \subset X_1 \subset \cdots \subset X_{n-1}
\]
en attachant à $X_{n-1}$ un espace $\widetilde{X}_n$ réunion disjointe de
$n$-boules fermées, par une application continue \emph{finie}
\[
  \partial\widetilde{X}_n \xrightarrow{\;q_n\;} X_{n-1}
\]
telle que l'image inverse de la filtration donnée de $X_{n-1}$ soit une
filtr.\ $n$-admissible de $\partial\widetilde{X}_n$, et\note{« $n$-admissible » tel qu'écrit ; par b) on attend $(n-1)$-admissible}

\page{36}
donne lieu à des applications \emph{étales}
\[
  q_n^{-1}(X_i) - q_n^{-1}(X_{i-1}) \longrightarrow X_i - X_{i-1} \qquad (0
  \leqslant i \leqslant n-1)
\]

Morphisme d'espaces à $n$-filtrations régulières
\[
  f : \underset{\substack{\| \\ X_n \\ \cup \\ \vdots}}{X} \longrightarrow
  \underset{\substack{\| \\ Y_n \\ \cup \\ \vdots}}{Y}
\]
a) $f$ fini

b) $f^{-1}(Y_i) = X_i$ \quad ($0 \leqslant i \leqslant n$)

c) \struck{\ill{}} $X_i - X_{i-1} \to Y_i - Y_{i-1}$ est étale (donc un
revêt.\ étale)

Catégorie $\mathcal{F}_n$ des espaces à $n$-filtrations régulières,
ss-catégorie pleine \struck{\ill{}} $\mathcal{F}_n^\circ$ des espaces
\uncertain{vérifiés} à $n$-filtration \emph{régulière}.
\note{il écrit « régulière » là où la page 35 disait « admissible » ; les
deux mots désignent visiblement la même notion}

\marginal{NB $X_{n-1} \to X_{n-1}$, $X_{n-1} \to \widetilde{X}_i \supset$
\ill{}, $X_{n-1}$ \ill{} \ill{} des foncteurs $\mathcal{F}_i$ \ill{}
$\mathcal{F}_n$ dans \ill{} $\mathcal{F}_i^\circ$ \ill{}
($\widetilde{X}_i$ \uncertain{vérifiant} $\partial\widetilde{X}_i \to
X_{i-1}$ \uncertain{homéomorphe}…)}
\note{note oblique dans la marge gauche, cernée d'un trait courbe, en
grande partie illisible}

$G_n$ groupe \emph{n}-cartographique (engendré par $\sigma_0, \sigma_1,
\ldots, \sigma_n$ avec $\sigma_i^2 = 1$ ($0 \leqslant i \leqslant n$),
$\sigma_i\sigma_j = \sigma_j\sigma_i$ si $n \geqslant j \geqslant i+2$, $i
\geqslant 0$ (condition \uncertain{si} $n \geqslant 2$))

Par récurrence sur $n$ on va définir des foncteurs

\begin{tikzcd}
  \mathcal{F}_n^\circ \arrow[r, "E_n^\circ"] \arrow[d, hook] & G_n\text{-ens}
  \arrow[d, "\text{oubli}"] \\
  \mathcal{F}_n \arrow[r, "E_n"] & G_{n-1}\text{-ens}
\end{tikzcd}

\qquad ($n \geqslant 1$)
\note{« $n \geqslant 1$ » est doublement souligné ; le signe à gauche
entre $\mathcal{F}_n^\circ$ et $\mathcal{F}_n$ est un $\cap$ (inclusion)}

\struck{(1) Cas $n = 1$ \quad $E_n(X_1) = \pi_0(\widetilde{X}_1^{\,\text{or}})
=$ ens des arcs de $X_1$, \uncertain{revêt.\ orienté} de
$\widetilde{X}_1$, $= \partial\widetilde{X}_1$ \ill{} ; $\sigma_0 =$
changement d'orientation ; $\sigma_1(a_0) = a'_0$ avec $q_0(a_0) =
q_0(a'_0)$, $a'_0 \neq a_0$, échange les pts des fibres de
$\partial\widetilde{X}_1 \to X_0$ ; $E_n^\circ(X_1) = \partial\widetilde{X}_1
\to \pi_0(\partial\widetilde{X}_1)$ ; $a_0 \in \partial\widetilde{X}_1$,
$a_1 \in \pi_0(\widetilde{X}_1)$, $a_0 \in a_1$, échangeant les pts des
fibres de …}
\note{tout le bas de la page, depuis « (1) Cas $n = 1$ » (le 1 cerclé),
est barré de trois longs traits obliques et surchargé ; on n'en donne que
ce qui se lit. À gauche, un petit diagramme $\widetilde{X}_1 \supset
\widetilde{X}_0 = \partial\widetilde{X}_1$ au-dessus de $X_1 \supset X_0$,
flèches $p_1$ et $q_0$, et le dessin d'une boucle fermée sur un point}

\page{37}
(avec diagrammes commutatifs).

(1) $n = 1$ \quad On pose pour $X_1 \in \mathrm{Ob}\, \mathcal{F}_n$
\[
  E_1(X_1) = \partial\widetilde{X}_1 \simeq \pi_0(\partial\widetilde{X}_1)
  \longrightarrow \pi_0(\widetilde{X}_1) \simeq \operatorname{arcs}(X_1)
\]
\note{« $\simeq$ arcs$(X_1)$ » est écrit sous $\pi_0(\widetilde{X}_1)$, avec
un signe $\simeq$ vertical ; le « 1 » initial est cerclé}
Les fibres de ce morphisme sont toutes de cardinal 2, d'où
\struck{\ill{}} \add{permutation} $\sigma_0$ commutant à ce morphisme,
sans pt fixe (« permutation caractéristique » de ce morphisme). Si $X_1$
est une 1-variété, l'appl.

\begin{tikzcd}
  E_1(X_1) = \partial\widetilde{X}_1 \arrow[d] \\
  X_0
\end{tikzcd}

a aussi ses fibres de card.\ 2, \struck{\ill{}} d'où perm.\ car.\ \add{de}
\uncertain{la} notée $\sigma_1$.

(2) Cas $n \geqslant 2$. On pose
\[
  E_n(X_n) = E^\circ_{n-1}(\partial\widetilde{X}_n) \quad (=
  E_{n-1}(\partial\widetilde{X}_n))
\]
Muni de $\sigma_0, \ldots, \sigma_{n-1}$ i.e.\ op.\ de $G_{n-1}$ par hyp.\
de récurrence. Fonctorialité évidente, \add{\uncertain{comme} composé de
foncteurs $\mathcal{F}_n \xrightarrow{\partial} \mathcal{F}_{n-1}^\circ
\xrightarrow{E^\circ_{n-1}} (G_{n-1}\text{-ens})$}. Si $X_n$ est une
$n$-variété, \struck{\ill{}} \add{l'application}

\begin{tikzcd}
  E_{n-1}(\partial\widetilde{X}_n) \arrow[d] \\
  E_{n-1}(X_{n-1})
\end{tikzcd}

(c'est un $G_{n-2}$-hom., i.e.\ commute à $\sigma_1, \ldots, \sigma_{n-2}$)
vérifie a des fibres de cardinal 2, l'hom.\ can.\ \struck{\ill{}} noté
$\sigma_n$, il commute à $\sigma_1, \ldots, \sigma_{n-2}$. OK.
Fonctorialité claire encore dans $\mathcal{F}_n^\circ$
\note{« vérifie » est écrit dans la marge gauche, devant une grande
parenthèse qui embrasse le diagramme et la ligne suivante ; « (2) » est
cerclé}

\page{38}
\emph{Th.\ 1} \add{Soit $n \geqslant 1$.} \struck{Le} foncteur $E_n^\circ$
\uncertain{passe} à la catégorie \uncertain{isotypique} \struck{\ill{}}
$\mathcal{F}_n^\circ$, et \struck{induit} un foncteur \emph{pl.\ fidèle}
\[
  \mathcal{F}_n^\circ \longrightarrow (G_n\text{-ens})
\]
(dont l'image essentielle est formée des $G_n$-ens.\ tels que les
$\sigma_i$ ($0 \leqslant i \leqslant n$) \add{y} opèrent
\struck{\uncertain{librement}} \struck{sans pts fixes}.
\note{« pl.\ fidèle » est souligné deux fois ; un trait vertical dans la
marge gauche accompagne l'énoncé ; la parenthèse n'est pas refermée}

Considérons $X_n \in \mathrm{Ob}\,\mathcal{F}_n$. On lui associe une
cascade d'ens.\ à opérateurs
\[
  \begin{array}{lll}
    R_n(X_n) \overset{\text{déf}}{=} E_n(X_n) & G_{n-1}\text{-ens} &
    \sigma_0 \cdots \sigma_{n-1} \\
    R_{n-1}(X_n) = E_{n-1}(X_{n-1}) & G_{n-2}\text{-ens} & \sigma_0 \cdots
    \sigma_{n-2} \\
    \cdots & & \\
    R_1(X_n) = E_1(X_1) & G_0\text{-ens} & \sigma_0 \\
    R_0(X_n) = E_0(X_0) \overset{\text{déf}}{=} X_0 & G_{-1}\text{-ens} &
    \text{i.e.\ ensemble}
  \end{array}
\]
\note{après $E_n(X_n)$, sur la première ligne, un trait qu'on pourrait
lire « $= 1$ » ; on ne le reproduit pas}

De plus on a \struck{\ill{}} \add{une} cascade d'applications
\[
  R_n(X_n) \to R_{n-1}(X_n) \to \cdots \to R_0(X_n)
\]
\struck{définies} \add{\ill{}}
\[
  \begin{array}{ccc}
    R_i(X_n) & \longrightarrow & R_{i-1}(X_n) \\
    \| & & \| \\
    E_i(X_i) & & E_{i-1}(X_{i-1}) \\
    \| & & \\
    E^\circ_{i-1}(\partial\widetilde{X}_i) = E_{i-1}(\partial\widetilde{X}_i)
    & &
  \end{array}
\]
(hom.\ de $G_{i-2}$-ens, commutant à $\sigma_0 \cdots \sigma_{i-2}$) se
déduit de $\partial\widetilde{X}_i \to X_{i-1}$ par application du
foncteur $E_{i-1}$.

\emph{Th 2} \add{Soit $n \geqslant 0$.} \struck{Soit} \uncertain{(Casca)}
la catégorie des « cascades d'ens.\ à op.\ » sous les groupes $G_{n-1},
G_{n-2}, \ldots, G_0$, $G_{-1} = \lbrace 1 \rbrace$, on a donc un foncteur
\note{un trait vertical dans la marge gauche commence avec l'énoncé du
Th 2}

\page{39}
\[
  \underset{\substack{| \\ \text{catégorie isotopique} \\ \text{déduite de }
  \mathcal{F}_n}}{\mathcal{F}_n} \xrightarrow{\;\varphi_n\;}
  (\mathrm{Casc})_n
\]
Ce foncteur est pl.\ fidèle.
\note{le trait vertical de la marge gauche se prolonge jusqu'ici : c'est
la fin de l'énoncé du Th 2}

Je \uncertain{faut} aussi trouver \struck{le} \add{un} foncteur
\struck{\ill{} \ill{} \ill{}.} \struck{\ill{}}
\[
  (G_n\text{-ens}) \xrightarrow{\;\psi_n\;} (\mathrm{Casc})_n
\]
rendant commutatif le diagramme

\begin{tikzcd}
  \mathcal{F}_n^\circ \arrow[r, hook] \arrow[d, hook] &
  (G_n\text{-ens}) \arrow[d, "\psi_n"] \\
  \mathcal{F}_n \arrow[r] & (\mathrm{Casc})_n
\end{tikzcd}

\note{l'objet en haut à droite, surchargé, se lit sans doute
« $G_n$-ens » ; la flèche horizontale supérieure et la verticale de
gauche partent d'un crochet (inclusions)}

Soit donc $X_n \in \mathrm{Ob}\,\mathcal{F}_n^\circ$

\emph{Th.\ 3}
\[
  \left\lbrace
  \begin{array}{ll}
    R_n(X_n) \simeq E_n^\circ(X_n) & G_{n-1}\text{-ens} \\
    R_{n-1}(X_n) \simeq E_n^\circ(X_n)/\sigma_n & G_{n-2}\text{-ens} \\
    R_{n-2}(X_n) \simeq E_n^\circ(X_n)/(\sigma_n, \sigma_{n-1}) &
    G_{n-3}\text{-ens} \\
    \vdots & \\
    R_1(X_n) \simeq E_n(X_n)/(\sigma_n, \ldots, \sigma_2) & \\
    R_0(X_n) \simeq E_n(X_n)/(\sigma_n, \ldots, \sigma_2, \sigma_1) &
  \end{array}
  \right.
\]
\marginal{(NB. on \uncertain{oublie} l'opération $\sigma_n$)}\note{écrit à
droite de la première ligne}
avec morphismes de transition définis par passage au quotient sans plus
\note{une petite coche dans la marge gauche, en face de cette ligne}

\page{40}
On voudrait trouver des foncteurs (adjoints ?) en sens inverse

\begin{tikzcd}
  \mathcal{F}_n^\circ \arrow[d, hook] & G_n\text{-ens} \arrow[l]
  \arrow[d, dashed] \\
  \mathcal{F}_n & (\mathrm{Casc})_n \arrow[l, "\psi_n"]
\end{tikzcd}

On définira $\psi_n$ par récurrence sur $n$.

(1) $n = 0$ \qquad $\psi_0(E_0) = E_0$ avec top.\ discr.

(2) Cas général.
\[
  \overbrace{E_n \to \underbrace{E_{n-1} \to \cdots \to E_1 \to
  E_0}_{\substack{\text{définit } X_{n-1} \supset X_{n-2} \supset \cdots \\
  \text{par hyp.\ de réc.} \\ X_{n-1} = \psi_{n-1}(C_{n-1})}}}^{C_n}
\]
\note{les chiffres (1), (2) sont cerclés ; sur la page, $C_{n-1}$ est
aussi marqué par une accolade au-dessus de $E_{n-1} \to \cdots \to E_0$}

On aura $\pi_0(\widetilde{X}_n) \simeq E_n / G_{n-1}$

$\partial\widetilde{X}_n$ est défini par \struck{\ill{}} ($E_n$ avec action
de $G_{n-1}$) par l'hyp.\ de récurrence
\[
  X_{n-1} = \psi_{n-1}(C_{n-1}), \qquad \partial\widetilde{X}_n =
  \psi_{n-1}(\rho_{n-1}(E_n))
\]
$\partial\widetilde{X}_n \to X_{n-1}$ déduit en appliquant $\psi_{n-1}$ à
\[
  \rho_{n-1}(E_n) \longrightarrow C_{n-1}
\]
\marginal{préciser ce point}
\note{$\rho_{n-1}(E_n)$ : la cascade déduite de $E_n$ muni de l'action de
$G_{n-1}$ ; la lettre $\rho$ est douteuse}

Donc on définira
\[
  \psi_n(C_n) = \text{multicône de } \psi_{n-1}(\rho_{n-1}(E_n) \to C_{n-1})
\]
\note{« multicône » (lecture douteuse) est écrit au-dessus d'un mot biffé
illisible}

--- --- Il faudrait donc expliciter en un

\emph{Th 4} Description de $\psi_n : (\mathrm{Casc})_n \to \mathcal{F}_n$,
et le cas échéant propriété d'adjonction.

\end{document}
